\chapter{Control identity, useful graphs and maintenance associations}\label{ch:identity}

\section{What changed in the 22 September revision}
The earlier demonstration could show empty positive-control charts when the loaded laboratory profile did not identify
the controls. An empty plot is not evidence that control measurements are absent. The corrected page distinguishes
configuration, eligibility and an incomplete statistical baseline. This revision keeps the 58-chart catalogue, removes
misleading pooled limits from the older control overview, and adds three descriptive graphs based on explicit control identity.

\begin{infobox}
The SC mapping is a \textbf{review-only profile}. It supplies candidate name rules; it does not certify the applicable
historical SOP, reference-material concentration, aliquot, or the interpretation of a sample-versus-LOD comparison.
The supplied procedures were used to define synthetic examples, not to execute laboratory work.
\end{infobox}

\section{Identity before statistics}
An ordinary specimen code is not a positive control simply because its vendor role is \emph{Unknown} or \emph{Calibrator}.
The page first matches a configured name rule and target. A negative matrix is expected not to amplify the screening target,
but its endogenous taxon target can amplify. That endogenous signal is a reference check, not automatically a fixed-input
internal positive control and not a blank.

For longitudinal control data the key retains the target, complete normalized material label, extraction date and any
aliquot or dilution suffix. The panel also requires the same instrument, recorded program fingerprint and firmware before
enabling statistical alarms. Matching by a shortened root is available in Review only as a descriptive view.
Different extraction batches are not pooled to estimate one control mean or SD.

The date parser accepts a calendar-valid day.month.four-digit-year label. It rejects impossible dates and malformed years,
and marks an extraction date after the run as inconsistent. It does not repair these labels by guessing. Unknown or invalid
dates are isolated by source in panel control keys. The original text remains visible and exportable.

\par\medskip\noindent{\small\textbf{Listing}~\texttt{sopControlIdentity()}\hfill\textcolor{gray}{HTML lines 4286--4308}}\par\nopagebreak
\lstinputlisting[style=vstyle,language=JavaScript,firstnumber=4286]{code/sopControlIdentity.js}


\par\medskip\noindent{\small\textbf{Listing}~\texttt{ccBaselineNeed()}\hfill\textcolor{gray}{HTML lines 4328--4328}}\par\nopagebreak
\lstinputlisting[style=vstyle,language=JavaScript,firstnumber=4328]{code/ccBaselineNeed.js}



\section{Which graphs are useful}
\begin{longtable}{@{}p{4cm}p{11cm}@{}}
\toprule View & Revised use\\\midrule\endhead
Positive-control Cq & One material, extraction batch, aliquot/dilution, target and instrument/program at a time. Non-detection
must be reviewed alongside the Cq of detected replicates.\\
Negative-control rate & Separate reaction blanks, extraction blanks, grinding/environment blanks and negative matrices.
Each stage has its own positive count and eligible denominator. Legacy pooled history remains explicitly labelled as stage unknown.\\
Endogenous reference & Inspect the biological reference separately; a fixed IC shift needs a configured reference Cq and a
matching IC target. Subtracting each run's own median would hide a common shift.\\
Root-name overview & Exploratory only. The older control graph retains the SOP range but no longer displays pooled mean/SD lines.\\
Thermal and mechanical plots & Compare comparable programs and instrument configurations. A trend is an investigation trigger,
not proof of component wear or a validated remaining-life estimate.\\
Operator funnel & Shows unadjusted event rates and workload. It does not adjust for extraction batch, assay, instrument or
case mix and cannot establish an operator effect.\\
\bottomrule
\end{longtable}

\section{New descriptive associations}
The three new plots use one synthetic control group per run, with detected and total replicate counts retained in their
CSV rows. Colours distinguish instrument, target, full control batch and program. Only finite detected Cq means enter the
scatter plots; missing Cq is not replaced by a cut-off.
\readingnote{These are descriptive associations. Read them with detection counts, program settings and maintenance records;
do not treat a slope as a failure probability.}

\fig{g_all_x_control_batches.png}{Control extraction batches over calendar time. The example deliberately changes from a
March extraction to a July extraction; these are separate series.}{control_batches}

\fig{g_all_x_control_age.png}{Control mean Cq against days since the labelled extraction. Invalid, missing or future
extraction dates are excluded from the plotted age values. The example includes an explicitly simulated age effect.}{control_age}

\fig{g_all_x_control_thermal.png}{Control mean Cq against measured peak cooling rate. Both change with time in this
synthetic example; the association does not distinguish ageing extract from changing cooling performance.}{control_thermal}

\par\medskip\noindent{\small\textbf{Listing}~\texttt{controlBatchRecords()}\hfill\textcolor{gray}{HTML lines 4335--4344}}\par\nopagebreak
\lstinputlisting[style=vstyle,language=JavaScript,firstnumber=4335]{code/controlBatchRecords.js}


\par\medskip\noindent{\small\textbf{Listing}~\texttt{controlBatchGraph()}\hfill\textcolor{gray}{HTML lines 4345--4361}}\par\nopagebreak
\lstinputlisting[style=vstyle,language=JavaScript,firstnumber=4345]{code/controlBatchGraph.js}



No correlation coefficient, significance test or causal score is assigned. Before estimating an equipment effect, match
material and extraction batch, target, dilution, program and analysis settings; inspect run order, service dates and operator
allocation. A prospective verification using a stable reference material is stronger evidence than a pooled correlation.
The current implementation groups by the recorded program string; it does not yet enforce a complete analysis-settings
fingerprint or model reagent lot, operator and extract age jointly.

\section{Baseline and alarm corrections}
Statistical limits and forecasts are suppressed until the configured Phase-I size is available (default 20, minimum configurable
size 5) in a single instrument/program/firmware cohort. Specification checks remain separate from statistical checks.
The default trend fit uses Phase I rather than continually refitting every new point. A zero-event p-chart baseline no longer
silently ignores a later positive event. It is still a simple binomial approximation; sparse-event uncertainty is not modelled
by an exact predictive interval.

The erroneous across-run $R_{4s}$ rule was removed. That rule needs appropriately standardized controls within a run;
adjacent run means are not a substitute. The remaining run rules are descriptive monitoring tools, not validated laboratory
acceptance criteria. Review uses exact names by default and a fixed initial baseline; root grouping cannot raise Westgard alarms.

\section{Demonstration and interpretation limits}
The regenerated set contains 57 synthetic runs and 325 synthetic history rows. Thirty screening runs include 24 runs of
the March control extraction and six of the July extraction. The short July batch intentionally illustrates an incomplete
baseline. Historical summary rows contain no dated positive-control metrics because their synthetic well-level batch identities
are not available. They continue to demonstrate thermal, optical, mechanical and process-stage histories.

The new plots show associations rather than equipment diagnoses. LED current, cover motion, cooling and operator differences
were planted independently by the generator. The maintenance suggestions in the catalogue are hypotheses to investigate.
Calibrated temperature verification, optical reference measurements and service records remain external evidence.


